\section{Grammar}
\label{sec:grammar}
	\begin{minted}{antlr}
	
	protocol : (preamble)? (varDef)? SEMICOLON (roleDef)+  CLPAR (protocolID ASSIGN blockStatement)* CRPAR ;
	
	protocolID : ID ;
	
	blockStatement : (statement)+;
	
	statement : branch | ifThenElse | end | internalAction | rec;
	
	branch : inputRole=role FROM outputRole+=role (COMMA outputRole+=role)* COLON LPAR (branchStat* | commStat )RPAR ;
	
	branchStat : (BRANCH SLPAR rateValues+=rate SRPAR updates DOT statement) ;
	
	commStat : (SLPAR rateValues+=rate SRPAR updates DOT statement) ;
	
	ifThenElse : IF cond AT role (COMMA cond AT role)* THEN CLPAR thenStat=statement CRPAR (ELSE CLPAR elseStat=statement CRPAR)?;
	
	rec : protocolID ;
	
	end : END (AT CLPAR roles+=role (COMMA roles+=role)* CRPAR)?;
	
	internalAction : SLPAR rate SRPAR message AT role DOT statement;
	
	updates : (SQLPAR (prec=actions) SQRPAR)? upds=message ;
	
	message : actions | first=DOUBLE_STRING loop | loop second=DOUBLE_STRING | loop*;
	
	actions : (action+=DOUBLE_STRING)? ( action+=DOUBLE_STRING)*  ;
	
	loop : FOREACH LPAR indexIteration=index op=(EQ | LE | GE | LEQ | GEQ | NEQ ) upperBound=index RPAR DOUBLE_STRING AT role;
	
	preamble : PREAMBLE (modelType SEMICOLON)? (variableDecl)* ENDPREAMBLE ;
	
	varDef : CHAR EQ INTEGER;
	
	modelType : DOUBLE_STRING;
	
	roleDef : role FROM (indexSpec)? roleSpec* SEMICOLON;
	
	roleSpec : (role (COLON (roleVar (COMMA roleVar)*))?);
	
	roleGroup : ID  ;
	
	roleIndex : ID SLPAR (index | (index BRANCH INTEGER)) SRPAR ;
	
	indexSpec : index IN SLPAR INTEGER DOTS upperBound=CHAR SRPAR ; 
	
	role : (roleGroup | roleIndex) ;
	
	roleVar : DOUBLE_STRING;
	
	variableDecl : DOUBLE_STRING;
	
	rate : DOUBLE_STRING ;
	
	cond : DOUBLE_STRING ; 
	
	index : CHAR ;
	
	SEMICOLON   : ';' ;
	COLON 		: ':' ;
	DOT 		: '.' ;
	DOTS 		: '...' ;
	COMMA		: ',' ;
	BRANCH 		: '+' ;
	FROM 		: '->';
	ASSIGN 		: ':=' ;
	NEQ 		: '!=' ;
	EQ 			: '=' ;
	LEQ			: '<=';
	GEQ			: '>=';
	LE			: '<';
	GE			: '>';
	UNDERSCORE 	: '_' ;
	STAR 		: '*' ;
	LPAR   		: '(' ;
	RPAR   		: ')' ;
	SLPAR   	: '[' ;
	SRPAR   	: ']' ;
	CLPAR  		: '{' ;
		CRPAR  		: '}' ;
	SQLPAR		: '<<' ;
	SQRPAR		: '>>' ;
	AT 			: '@' ;
	IF 			: 'if';
	THEN		: 'then';
	ELSE		: 'else';
	IFplus  	: 'IF' ;
	ELSEplus	: 'ELSE';
	END 		: 'END';
	AND			: '&&' ;
	PREAMBLE	: 'preamble';
	ENDPREAMBLE : 'endpreamble';
	CONST 		: 'const';
	FOREACH		: 'foreach';
	IN			: 'in' ;
	WS  		: [ \t\r\n\u00a0]+ -> skip ;
	DOUBLE_STRING : '"' ~('"')+ '"';
	
	//Numbers
	fragment DIGIT : '0'..'9';    
	INTEGER       : DIGIT+;
	
	//IDs
	CHAR  : 'a'..'z' |'A'..'Z' ;
	ID              :  (CHAR | UNDERSCORE)+ (CHAR | DIGIT | UNDERSCORE | STAR)* CHAR*;
	
	\end{minted}
\captionof{listing}{Grammar for ARIA}
\label{lst:full_grammar}
\newpage
\section{Vending Machine: \texttt{.sta} and \texttt{.tra} files}
\label{sec:tra-sta-vending-machine}
\begin{listing}[h]
	\centering
	\small % Schriftgröße etwas verkleinern, damit es gut nebeneinander passt
	
	% Linke Seite: State Vectors
	\begin{minipage}[t]{0.48\textwidth}
		\centering
		\textbf{State Vectors} \\
		\begin{minted}[fontsize=\footnotesize]{text}
			0: [1, 1, 0, 0, 0]
			1: [3, 1, 0, 0, 1]
			2: [4, 1, 0, 0, 2]
			3: [1, 1, 0, 1, 0]
			4: [1, 1, 1, 0, 0]
			5: [3, 1, 1, 0, 1]
			6: [4, 1, 1, 0, 2]
			7: [1, 1, 1, 1, 0]
			8: [3, 1, 1, 1, 1]
			9: [4, 1, 1, 1, 2]
			10: [3, 1, 0, 1, 1]
			11: [4, 1, 0, 1, 2]
		\end{minted}
	\end{minipage}
	\hfill % Schiebt die beiden minipages auseinander
	% Rechte Seite: Transitions
	\begin{minipage}[t]{0.48\textwidth}
		\centering
		\textbf{Transitions} \\
		\begin{minted}[fontsize=\footnotesize]{text}
			12 16
			0 1 1.0
			1 2 1.0
			2 3 0.9
			2 4 0.1
			3 10 1.0
			4 5 1.0
			5 6 1.0
			6 4 0.1
			6 7 0.9
			7 8 1.0
			8 9 1.0
			9 4 0.1
			9 7 0.9
			10 11 1.0
			11 3 0.9
			11 4 0.1
		\end{minted}
	\end{minipage}
	
	\caption{State vector mapping (\texttt{.sta}, left) and the resulting transition matrix (\texttt{.tra}, right) for the VendingMachine example}
	\label{lst:vending-states-matrix}
\end{listing}
\newpage
\section{The Dining Cryptographers Protocol}
\label{app:dining-cryptos}
\begin{listing}[h!]
	\begin{minted}{text}
		preamble
		"dtmc";
		"const int N = 3;"
		"const int p1 = 1;" "const int p2 = 2;" "const int p3 = 3;"
		endpreamble;
		
		crypt[i] -> i in [1...N] crypt[i] : "coin[i] = 0", "agree[i] = 0", "pay=0";
		
		{InitPay :=
			if "i = 1" @crypt[i] then {
				if "pay = 0" @crypt[i] then {
					crypt[i] -> crypt[i] : (+ ["0.25"] "pay'=0" . Crypto
					+ ["0.25"] "pay'=1" . Crypto
					+ ["0.25"] "pay'=2" . Crypto
					+ ["0.25"] "pay'=3" . Crypto)
				} else { Crypto }
			} else {
				if "pay > 0" @crypt[i] then { Crypto }}	
					
			Crypto := if "coin[i] = 0" @crypt[i]
			then { crypt[i] -> crypt[i] : 
			   (+ ["0.5"] "coin[i]' = 1" . Crypto2
				+ ["0.5"] "coin[i]' = 2" . Crypto2)}
							
			Crypto2 := if "coin[i]>0 && coin[i+1]>0" @crypt[i] then {
				if "coin[i] = coin[i+1]" @crypt[i] then {
					if "pay=p[i]" @crypt[i] then {
						["1.0"] "agree[i]'=0" @crypt[i] . Crypto3     } else {
						["1.0"] "agree[i]'=1" @crypt[i] . Crypto3 }}else {
					if "pay=p[i]" @crypt[i] then {
						["1.0"] "agree[i]'=1" @crypt[i] . Crypto3
					} else {
						["1.0"] "agree[i]'=0" @crypt[i] . Crypto3
			}}}	
		
			Crypto3 := END}		
	\end{minted}
	\caption{The Dining Cryptographers Protocol in the specified syntax}
	\label{lst:dining}
\end{listing}
